% Part: incompleteness
% Chapter: arithmetization-syntax
% Section: proofs-in-nd

\documentclass[../../../include/open-logic-section]{subfiles}

\begin{document}

\olfileid[fr]{inc}{art}{pnd}
\olsection{\usetoken{P}{derivation} en déduction naturelle}

\begin{explain}
Pour arithmétiser les !!{derivation}s, nous devons les représenter
par des nombres. Une !!{derivation} est un arbre de !!{formula}s
dont chaque inférence porte une ou deux étiquettes. Une
représentation récursive s'impose donc : nous représentons une
!!{derivation} par un uplet dont les composantes sont le nombre de
sous-!!{derivation}s immédiates aboutissant aux prémisses de la
dernière inférence, les représentations de ces sous-dérivations,
la !!{formula} finale, l'étiquette de décharge de la dernière
inférence et un nombre indiquant le type de cette dernière.
\end{explain}

\begin{defn}
  Si $\delta$ est !!a{derivation} en déduction naturelle, son numéro
  de G\"odel $\Gn{\delta}$ est défini inductivement comme suit :
  \begin{enumerate}
  \item Si $\delta$ ne comporte que l'hypothèse~$!A$, alors
    $\Gn{\delta}$ vaut $\tuple{0, \Gn{!A}, n}$. Le nombre~$n$ vaut
    $0$ si l'hypothèse ne porte pas d'étiquette ; sinon, il est son
    étiquette numérique strictement positive.\footnote{Précision éditoriale : une hypothèse étiquetée peut rester ouverte. Le code $0$ indique l'absence d'étiquette et n'autorise jamais une décharge ; la source le décrit simplement comme le cas non déchargé.}
  \item Si $\delta$ se termine par une inférence à zéro, une, deux
    ou trois prémisses, alors $\Gn{\delta}$ vaut respectivement
    \begin{align*}
      & \tuple{0, \Gn{!A}, n, k}, \\
      & \tuple{1, \Gn{\delta_1}, \Gn{!A}, n, k}, \\
      & \tuple{2, \Gn{\delta_1}, \Gn{\delta_2}, \Gn{!A}, n, k}, \text{ ou}\\
      & \tuple{3, \Gn{\delta_1}, \Gn{\delta_2}, \Gn{\delta_3}, \Gn{!A}, n,
        k}.
    \end{align*}
    Ici, $\delta_1$, $\delta_2$ et $\delta_3$ sont les
    sous-!!{derivation}s qui aboutissent aux prémisses de la dernière
    inférence de~$\delta$; $!A$ est la conclusion de cette inférence;
    $n$ est son étiquette de décharge ($0$ si aucune hypothèse n'est
    déchargée); enfin, $k$ est déterminé par le tableau suivant,
    selon la règle appliquée en dernier.

    \begin{tabular}{lccccccc}
      \text{Règle :} & \Intro{\land} & \Elim{\land} & \Intro{\lor} & \Elim{\lor} \\
      $k$: & 1 & 2 & 3 & 4  \\[1.5ex]
      \text{Règle :} &  \Intro{\lif} & \Elim{\lif} & \Intro{\lnot} & \Elim{\lnot} \\
      $k$: & 5 & 6 & 7 & 8 \\[1.5ex]
      \text{Règle :}  &  \FalseInt & \FalseCl & \Intro{\lforall} & \Elim{\lforall} \\
      $k$: & 9 & 10 & 11 & 12 \\[1.5ex]
      \text{Règle :} & \Intro{\lexists} & \Elim{\lexists} & \Intro{\eq} & \Elim{\eq} \\
      $k$: & 13 & 14 & 15 & 16
    \end{tabular}
  \end{enumerate}
\end{defn}

\begin{ex}
  Considérons la !!{derivation} très simple
  \begin{prooftree}
    \AxiomC{$\Discharge{!A \land !B}{1}$}
    \RightLabel{\Elim{\land}}
    \UnaryInfC{$!A$}
    \DischargeRule{\Intro{\lif}}{1}
    \UnaryInfC{$(!A \land !B) \lif !A$}
  \end{prooftree}
  Le numéro de G\"odel de l'hypothèse est
  $d_0 = \tuple{0, \Gn{(!A \land !B)}, 1}$. Le numéro de G\"odel de la
  !!{derivation} dont la conclusion résulte de~$\Elim{\land}$ est
  $d_1 = \tuple{1, d_0, \Gn{!A}, 0, 2}$ : $1$ indique que
  $\Elim{\land}$ a une seule prémisse; viennent ensuite $d_0$,
  le numéro de G\"odel de la conclusion~$!A$, le $0$ parce qu'aucune
  hypothèse n'est déchargée et le $2$ qui code $\Elim{\land}$. Le numéro de
  G\"odel de la !!{derivation} entière est donc
  $\tuple{1, d_1, \Gn{((!A \land !B) \lif !A)}, 1, 5}$,
  c'est-à-dire
  \[
  \tuple{1, \tuple{1, \tuple{0, \Gn{(!A \land !B)}, 1}, \Gn{!A}, 0, 2},
    \Gn{((!A \land !B) \lif !A)}, 1, 5}.
  \]
\end{ex}

\begin{explain}
Une fois cette représentation des !!{derivation}s fixée, nous
devons encore montrer que nous pouvons manipuler leurs numéros de
G\"odel par des opérations récursives primitives et exprimer ainsi
leurs propriétés et relations essentielles. Certaines opérations
sont simples : si $d$ est le numéro de G\"odel d'!!a{derivation},
$\fn{EndFmla}(d) = (d)_{(d)_0+1}$ donne le numéro de G\"odel de sa
!!{formula} finale, $\fn{DischargeLabel}(d) = (d)_{(d)_0 + 2}$ son
étiquette de décharge et $\fn{LastRule}(d) = (d)_{(d)_0 + 3}$ le
nombre indiquant le type de la dernière inférence. D'autres
propriétés sont bien plus difficiles à traiter; nous en
esquisserons au moins la méthode. Notre objectif est de montrer
que la relation « $\delta$ est !!a{derivation} de~$!A$ à partir
de~$\Gamma$ » est récursive primitive sur les numéros de G\"odel
de $\delta$ et de~$!A$.
\end{explain}

\begin{prop}
  Les relations suivantes sont récursives primitives :
  \begin{enumerate}
  \item $!A$ figure comme hypothèse dans~$\delta$ avec l'étiquette~$n$.
  \item Toutes les hypothèses de~$\delta$ portant l'étiquette~$n$
    sont de la forme~$!A$ : on peut donc !!{discharge} l'hypothèse~$!A$
    au moyen de cette étiquette dans~$\delta$.
  \end{enumerate}
\end{prop}

\begin{proof}
  Nous devons montrer que les relations correspondantes entre
  numéros de G\"odel de !!{formula}s et de !!{derivation}s sont
  récursives primitives.
  \begin{enumerate}
  \item Montrons que $\fn{Assum}(x, d, n)$, vraie si $x$ est le
    numéro de G\"odel d'une hypothèse de la !!{derivation} de numéro
    de G\"odel~$d$ portant l'étiquette~$n$, est récursive
    primitive. C'est le cas si la !!{derivation} de numéro de
    G\"odel~$\tuple{0, x, n}$ est une sous-!!{derivation} de
    celle codée par~$d$. Notre codage des !!{derivation}s est un cas
    particulier du codage des arbres introduit dans la
    \olref[cmp][rec][tre]{sec}. La fonction récursive primitive
    $\fn{SubtreeSeq}(d)$ donne donc une suite, éventuellement répétitive,
    des numéros de G\"odel de toutes les sous-!!{derivation}s
    de celle codée par~$d$. Nous pouvons alors définir
    \[
    \fn{Assum}(x, d, n) \defiff \bexists{i<\len{\fn{SubtreeSeq}(d)}}{(\fn{SubtreeSeq}(d))_i =
      \tuple{0, x, n}}.
    \]
  \item Montrons que $\fn{Discharge}(x, d, n)$, vraie si toutes les
    hypothèses de la !!{derivation} de numéro de G\"odel~$d$ portant
    l'étiquette~$n$ sont la !!{formula} de numéro de G\"odel~$x$,
    est récursive primitive. Cette relation est vraie si et seulement
    si $\bforall{y<d}{(\fn{Assum}(y, d, n) \lif y = x)}$.
  \end{enumerate}
\end{proof}

\begin{prop}
  \ollabel{prop:followsby}
  La propriété $\fn{Correct}(d)$, vraie si et seulement si la dernière
  inférence de la !!{derivation}~$\delta$ de numéro de G\"odel~$d$
  est correcte, est récursive primitive.
\end{prop}

\begin{proof}
  Nous devons montrer que, pour chaque règle d'inférence~$R$, la
  relation $\fn{FollowsBy}_R(d)$ est récursive primitive. Cette
  relation est vraie si et seulement si $d$ est le numéro de G\"odel
  d'!!a{derivation}~$\delta$ dont la !!{formula} finale découle, par
  une application correcte de~$R$, des sous-!!{derivation}s
  immédiates de~$\delta$.

  Le cas de la règle $\Intro{\land}$ est simple. Si $\delta$ se
  termine par une application correcte de cette règle, elle a la
  forme suivante :
  \begin{prooftree}
    \AxiomC{}
    \RightLabel{$\delta_1$}
    \DeduceC{$!A$}

    \AxiomC{}
    \RightLabel{$\delta_2$}
    \DeduceC{$!B$}

    \RightLabel{\Intro\land}
    \BinaryInfC{$!A \land !B$}
  \end{prooftree}
  Le numéro de G\"odel~$d$ de~$\delta$ est alors
  $\tuple{2, d_1, d_2, \Gn{(!A \land !B)}, 0, k}$, avec
  $\fn{EndFmla}(d_1) = \Gn{!A}$,
  $\fn{EndFmla}(d_2) = \Gn{!B}$, $n=0$ et $k=1$.
  Nous pouvons donc définir $\fn{FollowsBy}_{\Intro\land}(d)$ par
  \begin{multline*}
    (d)_0 = 2 \land \fn{DischargeLabel}(d) = 0 \land \fn{LastRule}(d) = 1 \land {}\\
    \fn{EndFmla}(d) = {}\\
    \Gn{(} \concat \fn{EndFmla}((d)_1) \concat \Gn{\land}
    \concat \fn{EndFmla}((d)_2) \concat \Gn{)}.
  \end{multline*}

  La règle $\Intro\eq$ fournit un autre exemple simple. Elle n'a
  aucune prémisse, donc $(d)_0 = 0$, comme dans le cas des
  hypothèses. Son étiquette de décharge est également nulle :
  $n=0$. En revanche, $!A$ doit être de la forme $\eq[t][t]$ pour
  un terme clos~$t$. On obtient ainsi la définition récursive
  primitive
  \begin{multline*}
    (d)_0 = 0 \land \fn{DischargeLabel}(d) = 0 \land {}\\
    \bexists{t<d}{(\fn{ClTerm}(t) \land
      \fn{EndFmla}(d) = {}\\
      \Gn{{\eq}(} \concat t \concat \Gn{,} \concat t \concat \Gn{)})}.
  \end{multline*}

  Pour un cas plus complexe, $\fn{FollowsBy}_{\Intro{\lif}}(d)$
  est vraie si et seulement si la !!{formula} finale de~$\delta$
  est de la forme $(!A \lif !B)$, si la !!{formula} finale de
  $\delta_1$ est~$!B$ et si toute hypothèse de~$\delta_1$ portant
  l'étiquette~$n$ est de la forme~$!A$. Nous pouvons l'exprimer
  au moyen de la relation récursive primitive suivante :
  \begin{multline*}
    (d)_0 = 1 \land {}\\
    \bexists{a<d}{((\fn{DischargeLabel}(d)=0\lor\fn{Discharge}(a, (d)_1, \fn{DischargeLabel}(d))) \land {}}\\
      \fn{EndFmla}(d) = (\Gn{(} \concat a \concat \Gn{\lif}
      \concat \fn{EndFmla}((d)_1) \concat \Gn{)}))
  \end{multline*}
  ($a$ représente ici le numéro de G\"odel de~$!A$.)

  Considérons enfin $\Intro{\lexists}$. La dernière inférence
  de~$\delta$ est correcte si et seulement s'il existe !!a{formula}~$!A$,
  un terme clos~$t$ et !!a{variable}~$x$ tels que
  $\Subst{!A}{t}{x}$ soit la !!{formula} finale de~$\delta_1$
  et $\lexists[x][!A]$ la conclusion de cette inférence. Ainsi,
  $\fn{FollowsBy}_{\Intro{\lexists}}(d)$ est vraie si et seulement si
  \begin{multline*}
    (d)_0 = 1 \land \fn{DischargeLabel}(d) = 0 \land {} \\
    \bexists{a < d}{\bexists{x<d}{\bexists{t<d}{
          (\fn{ClTerm}(t) \land \fn{Var}(x)  \land {}}}}\\
    \fn{Subst}(a,t,x) = \fn{EndFmla}((d)_1) \land
    \fn{EndFmla}(d) = (\Gn{\lexists} \concat x \concat a)).
  \end{multline*}

  Appelons $\fn{ShapeND}(d)$ la condition $\fn{SeqCode}(d)$ avec soit $(d)_0=0$ et $\len{d}=3$, soit $(d)_0\leq3$ et $\len{d}=(d)_0+4$. Nous définissons alors $\fn{Correct}(d)$ par
  \begin{multline*}
    \fn{ShapeND}(d)\land\fn{Sent}(\fn{EndFmla}(d)) \land \bigl[{}\\
    (\fn{LastRule}(d) = 1 \land
    \fn{FollowsBy}_{\Intro\land}(d)) \lor \dots \lor {}\\
    (\fn{LastRule}(d) = 16 \land \fn{FollowsBy}_{\Elim\eq}(d)) \lor {}\\
    \bexists{n<d}{\bexists{x<d}{(d = \tuple{0, x, n})}}\bigr].
  \end{multline*}
  La première ligne impose que la !!{formula} finale de la
  !!{derivation} codée par~$d$ soit
  un énoncé; la dernière couvre le cas où $d$ ne code qu'une
  hypothèse.
\end{proof}

\begin{prob}
  Définissez les propriétés suivantes sur le modèle de la
  \olref[inc][art][pnd]{prop:followsby} :
  \begin{enumerate}
  \item $\fn{FollowsBy}_{\Elim{\lif}}(d)$,
  \item $\fn{FollowsBy}_{\Elim{\eq}}(d)$,
  \item $\fn{FollowsBy}_{\Elim{\lor}}(d)$,
  \item $\fn{FollowsBy}_{\Intro{\lforall}}(d)$.
  \end{enumerate}
  Dans le dernier cas, vous devrez également montrer que l'on peut
  vérifier par un test récursif primitif si la dernière inférence de
  la !!{derivation} de numéro de G\"odel~$d$ satisfait la condition
  sur la variable propre : la variable propre~$a$ de l'inférence
  $\Intro{\lforall}$ ne doit apparaître ni dans la !!{formula}
  finale de la !!{derivation} codée par~$d$, ni dans une de ses
  hypothèses ouvertes. Vous pouvez
  utiliser à cette fin le prédicat récursif primitif
  $\fn{OpenAssum}$ de la \olref[inc][art][pnd]{prop:openassum}.
\end{prob}

\begin{prop}
  \ollabel{prop:deriv}
  La relation $\fn{Deriv}(d)$, vraie si $d$ est le numéro de G\"odel
  d'!!a{derivation} correcte~$\delta$, est récursive primitive.
\end{prop}

\begin{proof}
  Une !!{derivation}~$\delta$ est correcte si chacune de ses
  inférences applique correctement une règle, c'est-à-dire si
  chacune de ses sous-!!{derivation}s se termine par une inférence
  correcte. Ainsi, $\fn{Deriv}(d)$ est vraie si et seulement si
  \[
  \bforall{i<\len{\fn{SubtreeSeq}(d)}}{\fn{Correct}((\fn{SubtreeSeq}(d))_i)}
  \]
\end{proof}

\begin{prop}
  \ollabel{prop:openassum}
  La relation $\fn{OpenAssum}(z, d)$, vraie si $z$ est le numéro de
  G\"odel d'une hypothèse !!{undischarged}~$!A$ de la
  !!{derivation}~$\delta$
  de numéro de G\"odel~$d$, est récursive primitive.
\end{prop}

\begin{proof}
  Une occurrence d'hypothèse est !!{discharged} si elle porte
  l'étiquette~$n$ dans une sous-!!{derivation} de~$\delta$ qui se
  termine par une règle portant cette même étiquette de décharge.
  Ainsi, $!A$ est une hypothèse !!{undischarged} de~$\delta$
  lorsqu'au moins une de ses occurrences n'est pas déchargée.
  Cette distinction importe : $\delta$ peut contenir à la fois des
  occurrences déchargées et non déchargées de~$!A$.

  Considérons une suite $\delta_0$, \dots, $\delta_k$ telle que
  $\delta_0 = \delta$, que $\delta_k$ soit l'hypothèse
  $\Discharge{!A}{n}$ (pour un certain~$n$) et que
  $\delta_{i+1}$ soit une sous-!!{derivation} immédiate de
  $\delta_i$. S'il existe une telle suite sans qu'aucune des
  $\delta_i$ se termine par une inférence d'étiquette de
  décharge~$n$, alors $!A$ est une hypothèse !!{undischarged}
  de~$\delta$.

  La fonction récursive primitive $\fn{SubtreeSeq}(d)$ fournit
  une suite qui énumère les numéros de G\"odel de toutes les
  sous-!!{derivation}s de~$\delta$. Une suite de numéros de
  sous-!!{derivation}s suivant un chemin de la racine à une
  hypothèse en est une sous-suite. Pour ce qui suit, il suffit
  d'utiliser la condition récursive primitive
  $\fn{Subseq}(s, s')$ qui exige que chaque terme de~$s$ figure
  dans~$s'$ :
  $\bforall{i<\len{s}}{\lexists[j<\len{s'}][(s)_i = (s')_j]}$.
  Toute sous-suite vérifie cette condition; la condition de chemin
  sera imposée séparément. La relation « être une
  sous-!!{derivation} immédiate » est également récursive
  primitive : $\fn{Subderiv}(d, d')$ est vraie si et seulement si
  $\bexists{j<(d')_0}{d = (d')_{j+1}}$.
  Nous pouvons donc définir $\fn{OpenAssum}(z, d)$ par
  \begin{multline*}
    \bexists{s<\fn{sequenceBound}(d+1,d+1)}{(\fn{SeqCode}(s)\land0<\len{s}\land{}\\ \fn{Subseq}(s, \fn{SubtreeSeq}(d))
      \land (s)_0 = d \land {}} \\
    \bexists{n<d}{((s)_{\len{s} \tsub 1} = \tuple{0, z, n} \land {}}\\
      \bforall{i<(\len{s} \tsub 1)}{(\fn{Subderiv}((s)_{i+1}, (s)_i) \land {}}\\
      (n=0\lor\fn{DischargeLabel}((s)_i) \neq n)))).
  \end{multline*}
\end{proof}

\begin{prop}
  \ollabel{prop:prf-prim-rec}
  Supposons que $\Gamma$ soit un ensemble récursif primitif
  d'!!{sentence}s. Alors la relation $\Prf[\Gamma](x, y)$, qui
  exprime « $x$ est le code d'!!a{derivation}~$\delta$ de~$!A$ à
  partir d'hypothèses !!{undischarged}s appartenant à~$\Gamma$ et
  $y$ est le numéro de G\"odel de~$!A$ », est récursive primitive.
\end{prop}

\begin{proof}
  Supposons que l'appartenance $y \in \Gamma$ soit exprimée par le
  prédicat récursif primitif~$R_\Gamma(y)$. Nous devons montrer que
  $\Prf[\Gamma](x, y)$ est récursive primitive lorsqu'elle est vraie
  si et seulement si $y$ est le numéro de G\"odel d'un énoncé~$!A$
  et $x$ le code d'une !!{derivation} en déduction naturelle dont la
  !!{formula} finale est~$!A$ et dont toutes les hypothèses
  !!{undischarged}s appartiennent à~$\Gamma$.

  D'après la \olref{prop:deriv}, la propriété $\fn{Deriv}(x)$, vraie
  si et seulement si $x$ est le numéro de G\"odel d'!!a{derivation}
  correcte~$\delta$ en déduction naturelle, est récursive
  primitive. Nous pouvons donc définir $\Prf[\Gamma](x, y)$ par
  \begin{align*}
    \Prf[\Gamma](x, y) \defiff {}
    & \fn{Deriv}(x) \land \fn{EndFmla}(x) = y \land {} \\
    & \bforall{z < x}{(\fn{OpenAssum}(z, x) \lif R_\Gamma(z))}.
  \end{align*}
\end{proof}

\begin{explain}
Le test $\fn{OpenAssum}$ porte sur les occurrences d'hypothèses.
Le code $0$ d'étiquette représente une absence de décharge. Pour
une étiquette positive, on recherche un chemin de la racine
à une occurrence de l'hypothèse sur lequel aucune inférence ne
la décharge. Les codes diminuent strictement d'un nœud à son
fils. Partant de~$d$, un chemin contient donc au plus $d+1$
codes, chacun au plus égal à~$d$. La borne couvre le chemin
à un seul nœud quand la dérivation se réduit à une hypothèse.
\end{explain}

\end{document}
